% Generated draft: requires mathematical and source review.
\documentclass[11pt]{article}
\usepackage[margin=25mm]{geometry}
\usepackage{fontspec}
\setmainfont{Latin Modern Roman}
\usepackage{amsmath,amssymb,mathtools,unicode-math}
\setmathfont{Latin Modern Math}
\usepackage{enumitem,needspace,xurl,hyperref}
\setlength{\parindent}{0pt}
\setlength{\parskip}{5pt}
\setlength{\emergencystretch}{3em}
\newcommand{\N}{\mathbb{N}}
\newcommand{\Z}{\mathbb{Z}}
\newcommand{\Q}{\mathbb{Q}}
\newcommand{\R}{\mathbb{R}}
\newcommand{\C}{\mathbb{C}}
\newcommand{\F}{\mathbb{F}}
\newcommand{\A}{\mathbb{A}}
\newcommand{\PP}{\mathbb{P}}
\newcommand{\E}{\mathbb{E}}
\newcommand{\eps}{\varepsilon}
\newcommand{\dd}{\,\mathrm{d}}
\begin{document}
\section*{UnsolvedMath ID: 30005175}
\textbf{Source problem code:} catalogue entry
\subsection{Definitions and mathematical statement}
\textbf{Scope and unresolved hypotheses.} Let a Gushel--Mukai variety (GM variety) mean a variety in the sense of the cited Oberwolfach report. The supplied excerpt refers to a preceding intersection construction but does not reproduce it. Consequently, the precise class of varieties, including its smoothness and dimension restrictions, must be recovered from that report before this draft can serve as a complete standalone formulation. The following separates the supplied geometric definitions from conditional formulations of the conjectural targets.

\textbf{Generalized partners and duals over \(\C\).} The excerpt attaches to a GM variety \(X\) a Lagrangian data set \((V_6(X),V_5(X),A(X))\), where \(\dim V_6(X)=6\), \(V_5(X)\subset V_6(X)\) is a hyperplane, and \(A(X)\subset\bigwedge^3V_6(X)\) is Lagrangian for the alternating pairing
\[
\bigwedge^3V_6(X)\times\bigwedge^3V_6(X)\longrightarrow\bigwedge^6V_6(X),\qquad (u,v)\longmapsto u\wedge v.
\]
Let \(X,X'\) have dimensions \(n,n'\) with \(n\equiv n'\pmod 2\). They are generalized partners if there is an isomorphism \(f:V_6(X)\to V_6(X')\) such that \((\bigwedge^3f)(A(X))=A(X')\). They are generalized duals if there is an isomorphism \(f:V_6(X)\to V_6(X')^\vee\) such that
\[
(\bigwedge^3f)(A(X))=A(X')^\perp,
\]
where \(A(X')^\perp\subset\bigwedge^3(V_6(X')^\vee)\) is the annihilator under the natural dual pairing. No compatibility with the hyperplanes \(V_5\) is imposed by the supplied definition.

\textbf{Conjectural question.} The supplied edited statement asks whether the Hodge, Mumford--Tate, motivated Mumford--Tate, and Tate conjectures hold for GM varieties and their generalized partners or duals. It does not specify whether the intended additional target for a pair is a transfer principle, an isomorphism of motives, or simply the conjectures for each member individually.

For a smooth projective complex variety \(Y\) in the intended class, the usual rational Hodge conjecture asks, for every integer \(0\leq r\leq\dim Y\), whether
\[
\operatorname{im}\!\left(\operatorname{cl}_B:\operatorname{CH}^r(Y)\otimes_{\Z}\Q\longrightarrow H^{2r}(Y,\Q)\right)
=H^{2r}(Y,\Q)\cap H^{r,r}(Y).
\]
Here \(\operatorname{CH}^r\) is the group of codimension-\(r\) algebraic cycles modulo rational equivalence, and the intersection is taken inside complex cohomology. The record does not specify a restriction to particular codimensions or powers of \(Y\).

For clarity, the usual arithmetic formulations require additional data not supplied in the record: a finitely generated field \(k\) of characteristic zero, an embedding \(k\hookrightarrow\C\), and a smooth projective model \(Y/k\). With these explicitly auxiliary choices, write
\[
V_B^j=H^j(Y_{\C},\Q),\qquad V_\ell^j=H^j_{\mathrm{\acute et}}(Y_{\bar k},\Q_\ell),\qquad
G_\ell^j=\overline{\rho_{\ell,j}(\operatorname{Gal}(\bar k/k))}^{\,\mathrm{Zar}}.
\]
The usual Mumford--Tate equality, under Betti--\(\ell\)-adic comparison, is
\[
(G_\ell^j)^\circ=\operatorname{MT}(V_B^j)\otimes_{\Q}\Q_\ell.
\]
A full cohomological version asks for every degree \(j\) and every prime \(\ell\). The excerpt does not determine whether all degrees, only middle cohomology, or a primitive summand is intended.

For an Andr\'e motive \(M\) realizing the intended cohomology, let \(G_{\mathrm{mot}}(M)\) be the group of tensor automorphisms of the Betti realization on the tensor category generated by \(M\). The motivated Mumford--Tate target described in the excerpt identifies this group with the Mumford--Tate group and, after extension to \(\Q_\ell\), with the connected algebraic monodromy group. The excerpt also asserts that the relevant Andr\'e motives are abelian, meaning that they belong to the tensor category generated by motives of abelian varieties, allowing duals, Tate twists, and direct summands. The precise motive meant by ``their motives'' is not specified; abelianity and the group equalities must not be silently treated as the same assertion.

With the auxiliary arithmetic data above, a usual rational Tate formulation asks, for every finite extension \(k'/k\), every prime \(\ell\), and every \(0\leq r\leq\dim Y\), whether
\[
\operatorname{cl}_\ell:\operatorname{CH}^r(Y_{k'})\otimes_{\Z}\Q_\ell\longrightarrow
H^{2r}_{\mathrm{\acute et}}(Y_{\bar k},\Q_\ell(r))^{\operatorname{Gal}(\bar k/k')}
\]
is surjective. These field and quantifier choices are an explicit formulation template, not recovered hypotheses of the excerpt. The supplied literature assessment additionally mentions a \emph{generalized} Tate statement without defining it, so its target cannot be identified with this cycle-class assertion from the record alone.

\textbf{Supplied qualifications.} The excerpt states that generalized partners have isomorphic middle-degree Hodge structures up to a possible Tate twist. It attributes equivalence of Kuznetsov components for generalized partners or duals to Kuznetsov and Perry, reference [23]. Neither statement supplies a precise motive-level target for the edited question. The excerpt also describes the Tate conjecture as a formal consequence of the Hodge and Mumford--Tate conjectures; its intended hypotheses and realizations are not reproduced.

The registry labels the problem partially solved. Its supplied assessment credits Fu--Moonen with the motivated Mumford--Tate and generalized Tate statements for even-dimensional GM varieties and describes the odd-dimensional and broader claims as remaining open. This is supplied dated curation, not an independently verified theorem statement or current literature assessment.
\subsection{Short English statement}
The record asks whether several conjectures relating algebraic cycles, Hodge structures, Galois representations, and motives hold for Gushel--Mukai varieties, including generalized partners and duals. Their complex geometric relationship is defined, but the excerpt omits hypotheses needed for a unique arithmetic or motive-level statement. The supplied partially solved label and even-dimensional results require independent source review.
\subsection{Sources}
\begin{itemize}
\item Local registry: \texttt{data/unsolvedmath/problems.json}, permanent UnsolvedMath ID \(30005175\), exact source problem code \texttt{catalogue entry}. This draft uses the supplied complete registry record; no local file or external source was independently inspected.
\item Supplied source citation: \emph{Algebraic Geometry: Moduli Spaces, Birational Geometry and Derived Aspects} (2023), Oberwolfach Reports, pp.~1805--1864, DOI \texttt{10.4171/owr/2022/32}. Supplied source URL and literature-source link, labeled ``Oberwolfach Report 32/2022'': \url{https://doi.org/10.4171/owr/2022/32}.
\item Supplied literature source: Fu--Moonen, \emph{Algebraic cycles on Gushel--Mukai varieties}, \url{https://arxiv.org/abs/2207.01118}. The record attributes the even-dimensional motivated Mumford--Tate and generalized Tate results to these authors; their precise hypotheses and conclusions have not been independently verified here.
\item Literature assessment date: \texttt{2026-08-22}. This exact date records supplied dated curation, not independent verification. The supplied assessment calls the problem partially solved and describes odd-dimensional and broader claims as open.
\item Internal references in the excerpt: [1] for Andr\'e motives, [10] for GM data, [10] and [22] for generalized partners and duals, and [23] for the result attributed to Kuznetsov and Perry. Their bibliographic entries are absent from the supplied record. No individual proposer is supplied.
\end{itemize}
\end{document}
